% Évaluation : protocole, scénario nominal et robustesse (sous-sections de « Évaluation comparative », chapitre 4)
\subsection{Protocole}

Le protocole est celui du §\,7.3 du cahier des charges. Chaque politique est simulée sur les mêmes 30 années, de graines 9000 à 9029, qu'aucun réglage ni apprentissage n'a utilisées. Les nombres aléatoires communs font que deux politiques affrontent exactement les mêmes durées de vie et les mêmes aléas de délai : l'écart observé sur une graine ne doit rien au hasard du tirage. Les moyennes sont accompagnées de leur intervalle de confiance à 95~\%, et l'écart de coût de chaque politique à P1 est testé par un test de Wilcoxon apparié sur les 30 années.

\needspace{18\baselineskip}
\subsection{Scénario nominal}

\begin{tabularx}{\linewidth}{@{}l R P{3.1cm} r r r r r@{}}
\caption{Indicateurs moyens sur 30 années, scénario S1 nominal}\label{tab:evaluation}\\
\toprule
\hd{Pol.} & \hd{Coût (UR)} & \hd{IC 95~\%} & \hd{TRS} & \hd{Dispo.} & \hd{Ruptures} & \hd{Attente (h)} & \hd{$p$} \\
\midrule
\lignes{data/tableau_evaluation.tex}
\bottomrule
\end{tabularx}

La colonne $p$ donne la probabilité, sous l'hypothèse d'une absence d'écart avec P1, d'observer un écart de coût au moins aussi marqué. Les ruptures sont la part des demandes de pièces après panne que le stock n'a pas pu servir immédiatement ; l'attente est la durée moyenne d'attente de pièce par réparation, en heures. Le test de Wilcoxon conclut à une différence significative dès que celle-ci est réelle, mais il ne dit pas dans quel sens ; le test d'équivalence du tableau~\ref{tab:tost}, mené sur la même marge de $\pm 2$~\%, sépare donc les politiques qui s'écartent de P1 vers le mieux de celles qui s'en écartent vers le pire.

\needspace{13\baselineskip}
\begin{tabularx}{\linewidth}{@{}l R r l@{}}
\caption{Test d'équivalence (TOST) du coût annuel contre P1, marge $\pm 2$~\%, scénario S1 nominal}\label{tab:tost}\\
\toprule
\hd{Pol.} & \hd{Écart moyen (UR)} & \hd{Écart (\% de P1)} & \hd{Verdict} \\
\midrule
\lignes{data/tableau_tost.tex}
\bottomrule
\end{tabularx}

\intertitre{La prédiction réduit le coût, la surveillance d'état davantage encore}
P3, qui s'appuie sur le réseau de survie, coûte 49\,200 UR par an, 5,6~\% de moins que P1, et l'écart est significatif ($p < 10^{-6}$). Il ramène les pannes de 189 à 39 par an et porte la disponibilité à 99,25~\%, la meilleure des politiques réalistes. La maintenance conditionnelle P2 fait pourtant mieux, avec 47\,100 UR, soit 9,6~\% de moins que P1, sans aucun réglage. Elle ramène les pannes à 30 par an et, avec si peu de pannes, se contente des niveaux de stock de P0 : la valeur moyenne du stock est de 208 UR pour P2, contre 847 pour P3 et 1\,105 pour P1. Elle surveille aussi le taux de défauts, que le réseau de survie ignore : P3 laisse les composants s'user un peu plus, avec 2,4~\% de pièces défectueuses contre 2,2~\% pour P2. L'avantage de P2 tient donc moins à la qualité de sa détection qu'à ce qu'elle observe : la dérive du signal propre à chaque composant, que le simulateur relie directement à l'usure, est une information plus directe qu'un risque de panne estimé.

\intertitre{La fréquence de ruptures se lit avec le nombre de pannes}
P2 présente 23,9~\% de ruptures, mais sur 30 demandes seulement : cela représente 7,5 ruptures par an, contre 2,1 pour P1. Ces attentes, longues (110 heures en moyenne) mais rares, ne pèsent presque pas sur le TRS, proche de celui de P1 et de P3 (96,7 à 96,8~\%).

\intertitre{Les agents d'apprentissage par renforcement n'atteignent pas P1}
Les deux agents les plus proches de P1 sont ceux dont l'espace d'action se réduit à des seuils. P4, au niveau N1, est le plus proche : il coûte 52\,800 UR par an, seulement 1,3~\% de plus que P1, écart significatif mais d'ampleur marginale ($p = 0{,}019$), pour un TRS légèrement inférieur (96,44 contre 96,69~\%), un stock plus léger que P1 (631 UR contre 1\,105), davantage de ruptures (3,9 contre 1,1~\%) et un peu plus de pannes (203 contre 189 par an). Le DQN du niveau réduit, lui, coûte 54\,900 UR, 5,4~\% de plus que P1 : il fait ainsi à peine mieux que P0 (55\,400 UR), dont il partage le TRS (96,1 contre 96,1~\%), avec un stock à peine plus léger (172 UR contre 187) et un peu moins de pannes (177 contre 181 par an). Leur point commun est un espace d'action paramétrique : le DQN règle chaque jour un seuil d'âge commun à tous les composants et un niveau de stock commun, P4 un seuil d'âge et un niveau de stock par composant, c'est-à-dire une politique de la même famille que P1. Dans les deux cas, la maintenance préventive se déclenche sur l'âge rapporté à l'intervalle, non sur l'état observé du composant — là où P2 et P3 trouvent l'essentiel de leur avantage. Reste donc à savoir si ce vocabulaire interdit d'atteindre P1, ou si les agents l'ont simplement mal exploité.

\intertitre{Ce que le répertoire d'actions de chaque agent laisse atteindre}
Un agent ne peut pas descendre sous le meilleur coût de son propre répertoire. On appelle \emph{plancher de représentation} d'un espace d'action $A$ le coût de sa meilleure politique \emph{constante} — un réglage unique tenu toute l'année — évalué sur les mêmes trente années. L'écart de l'agent à ce plancher mesure ce qu'il n'a pas su tirer d'un répertoire qui le lui permettait ; l'écart du plancher à P1, ce que ce répertoire, sous sa forme la plus simple, laisse au-dessus de la référence. Les deux parts se somment exactement à l'écart observé. Pour l'espace réduit du DQN, les 36 actions sont énumérées, si bien que le plancher y est \emph{exact} ; pour l'espace N1 de P4 ($6^8 = 1\,679\,616$ actions), une montée de coordonnée depuis quatre départs en fournit un majorant réalisable — le coût retenu est celui d'une politique effectivement jouable.

\needspace{14\baselineskip}
\begin{tabularx}{\linewidth}{@{}l L R R R R@{}}
\caption{Décomposition de l'écart de coût de chaque agent à P1, scénario S1 nominal. Le plancher est le meilleur coût atteint par une politique constante de l'espace d'action ; l'écart d'apprentissage est ce que l'agent laisse entre lui et ce plancher, l'écart de représentation ce que le plancher laisse au-dessus de P1. Les pourcentages se rapportent au coût de P1.}\label{tab:plancher}\\
\toprule
\hd{Pol.} & \hd{Espace d'action} & \hd{Coût agent} & \hd{Plancher} & \hd{Écart d'apprentissage} & \hd{Écart de représentation} \\
\midrule
\lignes{data/tableau_plancher.tex}
\bottomrule
\end{tabularx}

Le tableau tranche, et de deux façons opposées. Le meilleur réglage fixe de l'espace réduit coûte 54\,279 UR, soit 4,1~\% au-dessus de P1 : le vocabulaire du DQN — seuil commun aux quatre composants, point de commande \emph{dérivé} du niveau de stock — ne peut pas rejoindre P1, et cette restriction explique 2\,143 UR des 2\,796 UR d'écart du DQN, l'apprentissage n'en portant que 652. La conclusion s'inverse pour l'espace N1 : la montée de coordonnée y trouve un réglage fixe à 50\,185 UR, \emph{1\,952 UR sous P1}. Le répertoire de P4 contient donc une politique qui satisfait le critère O4 : son écart observé de 662 UR se scinde en 2\,614 UR d'apprentissage et $-$1\,952 UR de représentation, la part de représentation étant \emph{négative} puisque son meilleur réglage fixe bat P1. Autrement dit, l'agent laisse inexploité l'avantage de 1\,952 UR que ses seuils lui offraient, et reste encore 2\,614 UR au-dessus du réglage fixe pourtant accessible dans son propre espace. Les politiques constantes formant un sous-ensemble de celles de l'espace, la part d'apprentissage est un \emph{minorant} et celle de représentation un \emph{majorant} : ces deux bornes vont dans le sens des conclusions ci-dessus.

Le verdict du critère O4 reste donc négatif, mais de justesse pour P4, comme le rappelle la conclusion. Son diagnostic, en revanche, n'est plus uniforme : le plafond des paramètres n'explique l'écart que pour le DQN, dont le vocabulaire est le plus pauvre ; pour P4, l'espace d'action \emph{contient} une solution qui bat P1, et c'est l'exploitation de cet espace qui fait défaut. P2, qui vaut 47\,100 UR sans apprentissage grâce à des décisions prises machine par machine, montre du reste qu'une telle solution n'a rien d'inaccessible : elle existe, les agents ne l'ont pas retrouvée.

Au niveau N2, où l'agent choisit lui-même la machine et le composant à intervenir, les agents restent loin du compte : le BDQ coûte 74\,800 UR, 43~\% de plus que P1, et P5 67\,000 UR, 28,6~\% de plus. P5 n'a pas appris à cibler la maintenance préventive : avec trois fois moins d'interventions préventives que P1 (188 contre 640 par an), il laisse survenir 339 pannes, contre 189, et son TRS tombe à 95,0~\%. Il remplace des composants, mais rarement ceux qui allaient tomber en panne. Le BDQ, qui entretient un peu plus (345 interventions par an) et subit autant de pannes (312), accumule surtout les commandes : son stock immobilisé atteint 5\,300 UR, près de cinq fois celui de P1. Avec le budget de calcul de ce projet, l'apprentissage par renforcement ne retrouve donc pas ce que les règles conçues à la main obtiennent ; la section~\ref{sec:conclusion} y revient.

\intertitre{Ce qui reste à gagner}
L'oracle, qui connaît la durée de vie réelle de chaque composant, coûte 42\,400 UR par an. L'écart de 10~\% avec P2 mesure la valeur d'une information parfaite sur l'état des composants ; il borne ce qu'un meilleur prédicteur pourrait encore apporter.

\begin{figure}[H]
\centering
\begin{tikzpicture}
\begin{axis}[rapport, height=7cm,
  xlabel={Coût total annuel (milliers d'UR)}, ylabel={TRS (\%)},
  xmajorgrids, ymajorgrids, restrict x to domain=40:66, enlarge y limits=0.12,
  xticklabel style={/pgf/number format/use comma, /pgf/number format/fixed, /pgf/number format/precision=1},
  yticklabel style={/pgf/number format/use comma, /pgf/number format/fixed, /pgf/number format/precision=1},
  nodes near coords, point meta=explicit symbolic,
  visualization depends on={value \thisrow{ancre} \as \ancre},
  every node near coord/.append style={font=\sffamily\scriptsize, text=primaire, anchor=\ancre, inner sep=3pt},
  ]
\addplot+[only marks, mark=*, mark size=2.4pt, color=primaire,
  error bars/.cd, x dir=both, x explicit, error bar style={primaire!45}]
  table[x=cout_k, y=trs, x error=ic_k, meta=politique, col sep=comma] {data/evaluation.csv};
\end{axis}
\end{tikzpicture}
\caption{Coût et TRS moyens de chaque politique ; les barres donnent l'intervalle de confiance du coût. Les politiques les plus favorables se trouvent en haut à gauche : coût faible, TRS élevé. P5 (67\,000 UR, TRS de 95,0~\%) et le BDQ (74\,800 UR, TRS de 95,6~\%) sortent du cadre.}
\label{fig:pareto}
\end{figure}

\subsection{Robustesse}

Les politiques réglées ou apprises sur le scénario nominal sont ensuite confrontées, sans nouveau réglage, à quatre situations qui s'en écartent : délais fournisseurs doublés, taux de panne augmenté de 30~\%, équipe réduite à un technicien, et flotte homogène S0. Le tableau~\ref{tab:robustesse} donne le coût moyen et sa variation par rapport au scénario nominal, le tableau~\ref{tab:robustesse-effets} reprend cette variation en la munissant de son intervalle de confiance et de sa taille d'effet, et la figure~\ref{fig:robustesse} compare ces variations d'une politique à l'autre. P3 s'appuie sur un réseau de survie dont l'entrée décrit la famille de chaque machine, et les quatre agents d'apprentissage par renforcement observent le risque prédit par ce même réseau : entraînés pour les deux familles du scénario nominal, ils ne peuvent pas scorer une flotte qui n'en compte qu'une, et leurs cellules restent vides en situation S0. Seules les politiques qui n'ont pas besoin de ce réseau, P0, P1, P2 et l'oracle, y restent évaluables.

\needspace{14\baselineskip}
\begin{tabularx}{\linewidth}{@{}l R R R R R@{}}
\caption{Coût moyen sur 30 années hors scénario nominal (milliers d'UR)}\label{tab:robustesse}\\
\toprule
\hd{Pol.} & \hd{Nominal} & \hd{Délais doublés} & \hd{Pannes $+$30~\%} & \hd{Un technicien} & \hd{Flotte S0} \\
\midrule
\lignes{data/tableau_robustesse.tex}
\bottomrule
\end{tabularx}

\needspace{16\baselineskip}
\begin{tabularx}{\linewidth}{@{}l R R R R@{}}
\caption{Variation appariée du coût par rapport au scénario nominal, intervalle de confiance à 95~\% et taille d'effet $d_z$. Une cellule vide signale les politiques dont l'observation dépend du réseau de survie, non évaluables sur une flotte à famille unique.}\label{tab:robustesse-effets}\\
\toprule
\hd{Pol.} & \hd{Délais doublés} & \hd{Pannes $+$30~\%} & \hd{Un technicien} & \hd{Flotte S0} \\
\midrule
\lignes{data/tableau_robustesse_effets.tex}
\bottomrule
\end{tabularx}

Les deux tableaux se lisent ensemble : le premier donne le niveau atteint, le second la confiance à accorder à l'écart. Chaque écart est mesuré sur les mêmes graines de simulation dans les deux situations, ce qui annule la part de variabilité commune aux politiques et ne laisse que l'effet propre à la situation. La colonne des délais doublés sépare nettement les familles : P1 et P3 subissent un surcoût de 5 à 6~\% seulement, quand P2, le DQN et P0 en accusent 25 à 26~\%. La réduction de l'équipe est le second discriminant : elle coûte 22~\% à P2 ($d_z = 6{,}1$) et 20~\% à P3, contre 10~\% à P4 et 2~\% à P5, dont le coût de main-d'œuvre est plus faible. Le BDQ est le plus sensible à cette réduction, avec 46,8~\% d'augmentation pour la plus faible taille d'effet du tableau ($d_z = 1{,}7$) : son écart-type d'un tirage à l'autre est grand, ce qui traduit une politique plus dispersée. Enfin, le passage à une flotte à famille unique abaisse le coût des quatre politiques évaluables de 9 à 10~%, mais l'effet reste modéré au regard de sa dispersion ($|d_z|$ de 1,0 à 1,6) : la disparition de la famille supplémentaire simplifie le réapprentissage comme la gestion, sans bouleverser l'ordre des politiques.

\begin{figure}[H]
\centering
\begin{tikzpicture}
\begin{axis}[rapport, ybar, bar width=4.2pt, width=\linewidth, height=7cm,
  ymin=-15, ymax=50,
  symbolic x coords={P0,P1,P2,P3,DQN,P4,P5,BDQ,Oracle}, xtick=data, enlarge x limits=0.07,
  xlabel={Politique}, ylabel={Variation du coût (\%)},
  yticklabel style={/pgf/number format/use comma, /pgf/number format/fixed, /pgf/number format/precision=0},
  legend style={at={(0.5,1.03)}, anchor=south, legend columns=4, fill=white, fill opacity=0.9, text opacity=1, font=\sffamily\footnotesize},
  legend cell align=left]
\addplot[fill=primaire, draw=none] table[x=politique, y=delais, col sep=comma] {data/robustesse.csv};
\addplot[fill=accent, draw=none] table[x=politique, y=pannes, col sep=comma] {data/robustesse.csv};
\addplot[fill=gris!60, draw=none] table[x=politique, y=technicien, col sep=comma] {data/robustesse.csv};
\addplot[fill=primaire!35, draw=none] table[x=politique, y=s0, col sep=comma] {data/robustesse.csv};
\draw[gris!70] (axis cs:P0,0) -- (axis cs:Oracle,0);
\legend{Délais doublés, Pannes $+$30~\%, Un technicien, Flotte S0}
\end{axis}
\end{tikzpicture}
\caption{Variation du coût annuel de chaque politique par rapport au scénario nominal, dans les quatre situations dégradées. Les barres les plus courtes désignent les politiques les plus robustes : P1 et P3 encaissent le doublement des délais en 5 à 6~\%, quand P2 y laisse 25~\%. Les colonnes S0 de P3 et des quatre agents appris (DQN, P4, P5, BDQ) sont absentes : le réseau de survie qui alimente leur observation n'est pas applicable à une flotte à famille unique.}
\label{fig:robustesse}
\end{figure}

\intertitre{Délais fournisseurs doublés}
C'est la situation qui sépare le plus les politiques. P1 et P3 ne voient leur coût croître que de 5 à 6~\% : P1 dispose de stocks élevés, et P3 dimensionne ses commandes sur la demande prévue pendant le délai, qui s'allonge avec lui. P2, dont le stock est bas, prend 25~\% et passe derrière elles. P3 devient la politique réaliste la moins chère, à 51\,700 UR, nettement devant P1 ($p < 10^{-5}$).

\intertitre{Taux de panne augmenté de 30~\%}
L'ordre du scénario nominal se maintient : P2 reste la moins chère, à 55\,000 UR, talonnée par P3 (55\,100 UR), toutes deux nettement devant P1 ($p < 10^{-5}$).

\intertitre{Un seul technicien}
Les politiques qui multiplient les interventions préventives souffrent le plus du manque de main-d'œuvre : P2 prend 22~\%, P3 20~\% et l'oracle lui-même 20~\%. P1 devient la politique réaliste la moins chère, à égalité statistique avec P2 ($p = 0{,}19$). Une politique fondée sur l'état des composants suppose donc une équipe capable d'absorber les interventions qu'elle déclenche ; les politiques étudiées se bornent à ne pas lancer plus d'interventions qu'il n'y a de techniciens libres le matin, sans planifier leurs remplacements pour lisser la charge.

\intertitre{Flotte homogène}
La flotte homogène S0 est la situation la moins exigeante : les politiques qui s'y appliquent y coûtent 9 à 10~\% de moins qu'en nominal. Comme P3, les quatre agents appris observent le risque prédit par le réseau de survie, dont l'entrée dépend du nombre de familles : ils ne sont pas applicables à une flotte à famille unique, et seule la comparaison des politiques sans réseau — P0, P1, P2 et l'oracle — y reste possible. Le classement du scénario nominal est conservé : P2 demeure la moins chère (42\,500 UR), devant P1 (47\,000 UR) et P0 (50\,100 UR), l'oracle restant en tête (38\,600 UR). Le rapport entre les politiques ne tient donc pas à la composition particulière de la flotte nominale.

\needspace{16\baselineskip}
\subsection{Généralisation du classement}

Le classement précédent vaut pour une composition précise de la flotte : cent machines réparties en deux familles, deux techniciens. L'item 12 du cahier des charges demande s'il résiste à un changement de composition. Cinq compositions ont donc été tirées par un plan de Dirichlet centré sur la composition nominale, avec une part de famille U de 0,10 à 0,88 et un effectif de un à trois techniciens. Le protocole est rejoué à l'identique sur les mêmes trente années, de graines 9\,000 à 9\,029. P1 et P3 ne peuvent pas être transférés en l'état : leurs seuils dépendent des lois de panne, donc de la part de chaque famille, et ils sont réglés composition par composition sur les graines 101 à 103. Les autres politiques — P0, P2, l'oracle et les quatre agents — sont conservées telles quelles ; réentraîner chaque agent par composition demanderait, au niveau N1, environ cinquante minutes par graine et par composition, hors du budget de calcul de ce projet. Le transfert des agents est donc une limite assumée de cette généralisation.

\begin{tabularx}{\linewidth}{@{}l r r r c R R P{2.4cm}@{}}
\caption{Stabilité du classement selon la composition de la flotte ; $\tau$ est le tau de Kendall du classement de la composition contre le classement nominal, et « ordre attendu » la conservation de la chaîne $P2 < P3 < P1 <$ agents du cahier des charges.}\label{tab:generalisation}\\
\toprule
\hd{Flotte} & \hd{Tech.} & \hd{Part U} & \hd{$\tau$} & \hd{Ordre attendu} & \hd{P1 (UR)} & \hd{P2 (UR)} & \hd{Meilleur agent} \\
\midrule
\lignes{data/tableau_generalisation.tex}
\bottomrule
\end{tabularx}

\begin{figure}[H]
\centering
\begin{tikzpicture}
\begin{axis}[rapport, ybar, bar width=14pt, width=0.62\linewidth, height=5.6cm,
  ymin=0, ymax=1.08, ytick={0,0.25,0.5,0.75,1},
  symbolic x coords={C0,C1,C2,C3,C4}, xtick=data, enlarge x limits=0.16,
  xlabel={Composition de flotte}, ylabel={$\tau$ de Kendall contre le nominal},
  yticklabel style={/pgf/number format/use comma, /pgf/number format/fixed, /pgf/number format/precision=2}]
\addplot[fill=primaire, draw=none] table[x=composition, y=tau, col sep=comma] {data/generalisation.csv};
\draw[accent, densely dashed, thick] (axis cs:C0,1) -- (axis cs:C4,1);
\end{axis}
\end{tikzpicture}
\caption{Concordance du classement de chaque composition tirée avec le classement nominal. La ligne en pointillés marque un classement identique ($\tau = 1$) : les barres s'en approchent toutes, sauf la composition C0, à un seul technicien.}
\label{fig:generalisation}
\end{figure}

\intertitre{Un classement qui résiste, à une exception près}
Le classement se reproduit presque exactement d'une composition à l'autre. Le tau de Kendall vaut 1,000 pour deux compositions (C2 et C4), 0,944 pour deux autres (C1 et C3), et 0,667 pour la seule composition à un technicien (C0). Deux compositions renversent l'ordre attendu par le cahier des charges. En C0, à équipe réduite, P1 — dont le stock élevé absorbe les délais — passe devant P3 et P2 et devient, après l'oracle, la politique réaliste la moins chère. En C3, à trois techniciens et forte part de la famille soumise à l'usure, P4 passe devant P1 (56\,419 UR contre 56\,676 UR) : c'est le seul cas, sur les six compositions, où un agent devance la meilleure politique à seuils. Partout ailleurs, P2 demeure la moins chère des politiques sans oracle et les quatre agents restent derrière P1. Le meilleur agent est P4 dans les cinq compositions comme en nominal : le résultat du scénario nominal — les agents n'atteignent P1 que dans une composition particulière, et jamais de façon robuste — ne dépend donc pas de la composition.

\intertitre{Le niveau du coût suit la composition, l'ordre non}
Le coût annuel de P1 varie du simple au double selon la composition, de 33\,900 UR pour une flotte presque exclusivement calendaire à 65\,900 UR pour une flotte presque exclusivement soumise à l'usure, et les autres politiques se déplacent avec lui. Ce qui change d'une composition à l'autre est donc le niveau du coût, non la hiérarchie des politiques. L'écart établi au scénario nominal n'est pas un artefact du jeu d'essai.
